\subsection{Synthesis: Architectural Split Choice}
\label{sec:analysis-split-synthesis}

The three preceding subsections establish the local architectural
split-choice stage of the thesis. The two-dimensional explanatory
account developed in \cref{sec:analysis-rq13} --- activation-transfer
size at the coarse level and compute placement at the boundary at the
refined level --- is not restated here; the role of this subsection is
to make the section's contribution to the rest of the thesis explicit
and to set the reading rule against which the Kubernetes and Azure
stages are interpreted.

The local stage contributes two things forward. First, an empirically
selected carry-forward set --- \texttt{split\_after\_layer3\_block0}
(rule-selected main) and \texttt{split\_after\_layer2} (reference) ---
which fixes the boundary configurations that the subsequent stages
re-measure rather than re-search. Second, a joint suitability criterion
in which activation-transfer size and compute placement are
complementary, not competing, dimensions; this criterion rejected the
two coarse endpoints (\texttt{split\_after\_layer1} on activation
expansion, \texttt{split\_after\_layer4} on compute degeneracy) and
identified the mid-to-late band as the only region in which both
dimensions are simultaneously tolerable under the conditions measured
here. Both contributions are properties of the CPU-only batch-one
localhost-gRPC environment defined in
\cref{tab:rq11-config,tab:rq12-config} and remain bounded by the
single-host, single-build, fixed-seed scope of those runs.

What the local stage does \emph{not} establish is whether the joint
criterion continues to identify the same region once Kubernetes
networking, sidecar mediation, or managed-cloud node variability is
added to the path. This is the empirical question that the Kubernetes
(Stage~K) and Azure (Stage~C) stages are designed to answer. The reading
rule against which their results are interpreted in
\cref{sec:analysis-deployment} is trend transfer rather than
millisecond identity: the comparison preserves the ranking and
direction of the local explanatory account where the underlying
mechanism is plausibly stable across environments, and treats absolute
latency shifts as legitimate consequences of the new
substrate~\cite{benchmarking_sc15,noise_in_clouds}.
